% AUTO-GENERATED by src/eda.py - do not edit by hand
\providecommand{\NRawRows}{100,775}
\providecommand{\NAfterItemDedup}{63,773}
\providecommand{\NAfterResidential}{57,796}
\providecommand{\NAfterValidity}{57,712}
\providecommand{\NFinal}{57,480}
\providecommand{\NTrain}{40,236}
\providecommand{\NVal}{8,622}
\providecommand{\NTest}{8,622}
\providecommand{\NFeatures}{27}
\providecommand{\NFlats}{56,391}
\providecommand{\NSharedFlats}{0}
\providecommand{\TierThreshold}{208,000}
\providecommand{\PremiumTrain}{0.500}
\providecommand{\PremiumVal}{0.500}
\providecommand{\PremiumTest}{0.500}
\providecommand{\LeakTotalPriceCorr}{1.0000}
\providecommand{\LeakUnitPriceCorr}{1.0000}

\begin{table*}[t]
\centering\small
\caption{Task A (price regression), test set. $\pm$: bootstrap std; best per column in bold.}\label{tab:results_a}
\begin{tabular}{lrrrrrr}
\toprule
Model & RMSE (log) & MAE (log) & $R^2$ (log) & $R^2$ (AZN) & MAE (AZN) & Fit (s) \\
\midrule
Mean predictor & 0.629 $\pm$ 0.006 & 0.488 $\pm$ 0.004 & -0.000 $\pm$ 0.000 & -0.049 $\pm$ 0.004 & 127,850 $\pm$ 2,295 & 0.00 \\
Ridge, normal eq. (ours) & \textbf{0.226 $\pm$ 0.003} & \textbf{0.157 $\pm$ 0.002} & \textbf{0.870 $\pm$ 0.004} & 0.705 $\pm$ 0.031 & \textbf{48,280 $\pm$ 1,405} & 0.01 \\
Decision tree (ours) & 0.239 $\pm$ 0.004 & 0.163 $\pm$ 0.002 & 0.856 $\pm$ 0.004 & 0.705 $\pm$ 0.034 & 49,516 $\pm$ 1,413 & 12.46 \\
sklearn DecisionTreeRegressor & 0.238 $\pm$ 0.004 & 0.162 $\pm$ 0.002 & 0.856 $\pm$ 0.004 & \textbf{0.705 $\pm$ 0.034} & 49,353 $\pm$ 1,416 & 0.91 \\
\bottomrule
\end{tabular}
\end{table*}

\begin{table*}[t]
\centering\small
\caption{Task B (price tier), test set. $\pm$: bootstrap std; best per column in bold.}\label{tab:results_b}
\begin{tabular}{lrrrrrrr}
\toprule
Model & Acc. & Prec. & Rec. & F1 & ROC-AUC & PR-AUC & Fit (s) \\
\midrule
Majority class & 0.500 $\pm$ 0.006 & 0.000 $\pm$ 0.000 & 0.000 $\pm$ 0.000 & 0.000 $\pm$ 0.000 & 0.500 $\pm$ 0.000 & 0.500 $\pm$ 0.006 & 0.00 \\
Decision tree (ours) & 0.914 $\pm$ 0.003 & 0.916 $\pm$ 0.004 & 0.912 $\pm$ 0.004 & 0.914 $\pm$ 0.003 & 0.972 $\pm$ 0.002 & 0.968 $\pm$ 0.002 & 65.39 \\
sklearn DecisionTreeClassifier & 0.914 $\pm$ 0.003 & 0.916 $\pm$ 0.004 & 0.912 $\pm$ 0.004 & 0.914 $\pm$ 0.003 & 0.972 $\pm$ 0.002 & 0.969 $\pm$ 0.002 & 0.66 \\
Pegasos SVM (ours) & 0.906 $\pm$ 0.003 & 0.889 $\pm$ 0.005 & 0.928 $\pm$ 0.004 & 0.908 $\pm$ 0.003 & 0.968 $\pm$ 0.002 & 0.967 $\pm$ 0.002 & 1.05 \\
sklearn SVC & \textbf{0.924 $\pm$ 0.003} & \textbf{0.920 $\pm$ 0.004} & \textbf{0.929 $\pm$ 0.004} & \textbf{0.924 $\pm$ 0.003} & \textbf{0.981 $\pm$ 0.001} & \textbf{0.981 $\pm$ 0.001} & 306.22 \\
sklearn SGD (hinge) & 0.912 $\pm$ 0.003 & 0.902 $\pm$ 0.004 & 0.924 $\pm$ 0.004 & 0.913 $\pm$ 0.003 & 0.969 $\pm$ 0.002 & 0.967 $\pm$ 0.002 & 0.07 \\
\bottomrule
\end{tabular}
\end{table*}

\begin{table*}[t]
\centering\small
\caption{Task A, nested cross-validation on train+val (mean $\pm$ std over folds).}\label{tab:cv_a}
\begin{tabular}{lrrr}
\toprule
Model & RMSE (log) & MAE (log) & $R^2$ (log) \\
\midrule
Mean predictor & 0.637 $\pm$ 0.004 & 0.492 $\pm$ 0.002 & -0.000 $\pm$ 0.000 \\
Ridge, normal eq. (ours) & \textbf{0.226 $\pm$ 0.003} & \textbf{0.156 $\pm$ 0.002} & \textbf{0.874 $\pm$ 0.004} \\
Decision tree (ours) & 0.236 $\pm$ 0.003 & 0.160 $\pm$ 0.002 & 0.863 $\pm$ 0.002 \\
sklearn DecisionTreeRegressor & 0.236 $\pm$ 0.003 & 0.160 $\pm$ 0.002 & 0.863 $\pm$ 0.002 \\
\bottomrule
\end{tabular}
\end{table*}

\begin{table*}[t]
\centering\small
\caption{Task B, nested cross-validation on train+val (mean $\pm$ std over folds).}\label{tab:cv_b}
\begin{tabular}{lrrr}
\toprule
Model & Acc. & F1 & ROC-AUC \\
\midrule
Majority class & 0.501 $\pm$ 0.006 & 0.000 $\pm$ 0.000 & 0.500 $\pm$ 0.000 \\
Decision tree (ours) & 0.913 $\pm$ 0.004 & 0.913 $\pm$ 0.004 & 0.972 $\pm$ 0.001 \\
sklearn DecisionTreeClassifier & 0.913 $\pm$ 0.004 & 0.913 $\pm$ 0.004 & 0.971 $\pm$ 0.001 \\
Pegasos SVM (ours) & 0.910 $\pm$ 0.003 & 0.910 $\pm$ 0.004 & 0.969 $\pm$ 0.002 \\
sklearn SVC & \textbf{0.925 $\pm$ 0.003} & \textbf{0.925 $\pm$ 0.003} & \textbf{0.980 $\pm$ 0.001} \\
sklearn SGD (hinge) & 0.912 $\pm$ 0.003 & 0.912 $\pm$ 0.003 & 0.970 $\pm$ 0.002 \\
\bottomrule
\end{tabular}
\end{table*}

\begin{figure}[t]\centering
\includegraphics[width=\columnwidth]{figures/roc_pr_taskB.png}
\caption{Task B ROC and precision--recall curves on the test set.}\label{fig:roc}
\end{figure}
\begin{figure}[t]\centering
\includegraphics[width=\columnwidth]{figures/pred_vs_true_taskA.png}
\caption{Task A: predicted vs true price (best model) and log residuals.}\label{fig:pred}
\end{figure}
% Emil:
Error analysis. Errors are not uniform. Mid-priced listings (100–400k AZN) have a median
absolute percentage error of about 10%, but the model regresses toward the mean: it
overestimates cheap listings (+15% bias below 100k AZN) and underestimates luxury
listings (−27% above 800k AZN, mdape 32%). Houses are the hardest category
(mdape 22.7% vs. about 10% for flats), and all 20 largest absolute errors are
underpredictions, 18 of them large houses (mean area about 613 m$^2$) in affluent areas
such as Badamdar, Mərdəkan and Bilgəh. Price depends there on land, view and construction
quality, which our features do not capture (the free-text description was not used), and
a log-scale linear/tree model compresses the extreme upper tail. Future work: land-area
interactions, text features, quantile or two-stage models for the luxury segment.
% report/figures/error_by_*.csv and results.json ("error_analysis_taskA").
